\documentclass[12pt, titlepage]{article}
\usepackage{amsmath}
\usepackage[round]{natbib}
\usepackage{amssymb}
\usepackage{graphicx}
\usepackage{hyperref}
\usepackage{longtable}
\usepackage{tabularx}
\usepackage{float}
\usepackage{booktabs}
\usepackage{multirow}
\usepackage{fullpage}
\hypersetup{
bookmarks=true,
colorlinks=true,
linkcolor=red,
citecolor=blue,
filecolor=magenta,
urlcolor=cyan
}
\usepackage{array}
\input{../../Comments.text}
\input{../../Common.text}

\begin{document}
\sloppy

\title{Module Interface Specification for \progname{}\\ \large Design Document Revision $-1$}
\author{\authname}
\date{\today}

\maketitle

\pagenumbering{roman}

\section{Revision History}

\begin{tabularx}{\textwidth}{p{3cm} p{3cm} X}
\toprule
\textbf{Date} & \textbf{Version} & \textbf{Notes}\\
\midrule
Nov 16, 2026 & Rev $-1$ & Module interface syntax: access programs, state variables, environment variables and exceptions. Access routine semantics are deferred to Rev 0.\\
\bottomrule
\end{tabularx}

\newpage
\section{Symbols, Abbreviations and Acronyms}

See the SRS \citep{TandemSRS}
for a complete list of symbols, abbreviations and acronyms. The additional terms
used in this document are MG (Module Guide), MIS (Module Interface
Specification) and ADT (abstract data type).

\newpage
\tableofcontents
\newpage

\pagenumbering{arabic}

\section{Introduction}

This document gives the interfaces of the modules in the
Module Guide \citep{TandemMG}, in the same order and grouping. Revision $-1$ gives syntax only; access routine
semantics are deferred to Revision 0.

\section{Notation and Types}

The MIS structure and notation follow \citet{HoffmanAndStrooper1995}. The
primitive data types are:

\begin{center}
\renewcommand{\arraystretch}{1.2}
\begin{tabular}{l l p{7.5cm}}
\toprule
\textbf{Data Type} & \textbf{Notation} & \textbf{Description}\\
\midrule
character & char & a single symbol or digit\\
integer & $\mathbb{Z}$ & a number without a fractional component in $(-\infty, \infty)$\\
natural number & $\mathbb{N}$ & a number without a fractional component in $[0, \infty)$\\
real & $\mathbb{R}$ & any number in $(-\infty, \infty)$\\
boolean & $\mathbb{B}$ & $\{\mathit{true}, \mathit{false}\}$\\
\bottomrule
\end{tabular}
\end{center}

The derived types are listed below. A tuple is written with its field names and
types; $A \to B$ is a function (map) from $A$ to $B$, and $\mathcal{P}\,A$ is a
set of elements of $A$.

\subsubsection*{Text, models and devices}

\begin{description}
\item[String:] a sequence of characters.
\item[Token:] an integer token ID of one tokenizer; \texttt{Seq of Token} is a
  sequence of token IDs.
\item[ModelName:] one of the answer models: $\{\mathit{Qwen}$, $\mathit{SmolLM}\}$
  in the proof of concept.
\item[Revision:] a pinned model or tokenizer revision, a String.
\item[Model:] an answer model loaded by Model Runtime (\hyperref[mModelRun]{M12}).
\item[ClassifierModel:] the classifier loaded by Model Selection
  (\hyperref[mSelect]{M5}).
\item[DeviceName:] one of $\{\mathit{cpu}$, $\mathit{cuda}$, $\mathit{mps}\}$.
\item[DType:] one of $\{\mathit{float32}$, $\mathit{float16}$, $\mathit{bfloat16}\}$, a
  number format.
\item[DeviceHandle:] a tuple $(\mathit{device}: \mathit{DeviceName}$, $\mathit{dtype}:
  \mathit{DType})$, a device reserved by Accelerator Interface
  (\hyperref[mAccel]{M1}).
\item[File:] a file on disk.
\end{description}

\subsubsection*{Caches}

\begin{description}
\item[Cache:] an instance of the Cache Store ADT (\hyperref[mCache]{M13}).
\item[Keys:] Seq of Seq of Seq of (Seq of $\mathbb{R}$): one vector per layer, head
  and token position.
\item[KVData:] a tuple $(\mathit{keys}: \mathit{Keys}$, $\mathit{values}: \mathit{Keys})$.
\item[Provenance:] a tuple $(\mathit{model}: \mathit{ModelName}$, $\mathit{tokens}:
  \text{Seq of Token}$, $\mathit{complete}: \mathbb{B})$, the tags of a Cache.
\item[CheckpointRef:] a tuple $(\mathit{name}: \mathit{String}$, $\mathit{direction}:
  \mathit{ModelName} \times \mathit{ModelName}$, $\mathit{revision}: \mathit{Revision}$, $\mathit{identifier}: \mathit{String})$.
\item[Map:] the translator loaded from one checkpoint, a function from KVData to
  KVData.
\item[CacheRecord:] a tuple $(\mathit{source}: \mathit{ModelName}$, $\mathit{target}:
  \mathit{ModelName}$, $\mathit{predictor}: \mathit{String}$, $\mathit{targetRevision}:
  \mathit{Revision}$, $\mathit{alignment}: \mathit{String}$, $\mathit{checksum}: \mathit{String})$.
\item[Reason:] one of $\{\mathit{prefixMismatch}$, $\mathit{wrongModel}$, $\mathit{wrongRevision}$, $\mathit{shapeMismatch}$, $\mathit{dtypeMismatch}$, $\mathit{nonFinite}$, $\mathit{noPermission}$, $\mathit{checksumMismatch}\}$.
\item[Verdict:] either $\mathit{accepted}$ or $\mathit{rejected}(\mathit{Reason})$.
\end{description}

\subsubsection*{Requests and conversations}

\begin{description}
\item[RequestID:] a unique identifier of a request, a String.
\item[ReuseDomain:] a host-assigned group ID that grants sharing permission, a
  String. Only the host sets it; a request cannot set or change it.
\item[Settings:] a tuple $(\mathit{maxNewTokens}: \mathbb{N})$.
\item[Request:] a tuple $(\mathit{id}: \mathit{RequestID}$, $\mathit{prompt}:
  \mathit{String}$, $\mathit{conversation}: \mathit{String}$, $\mathit{current}:
  \mathit{ModelName}$, $\mathit{domain}: \mathit{ReuseDomain} \cup \{\bot\}$, $\mathit{settings}:
  \mathit{Settings})$.
\item[ErrorCode:] one of $\{\mathit{missingField}$, $\mathit{emptyPrompt}$, $\mathit{promptTooLong}$, $\mathit{invalidSettings}\}$.
\item[Answer:] a String of generated text.
\item[Turn:] a tuple $(\mathit{role}: \{\mathit{user}$, $\mathit{assistant}\}$, $\mathit{text}: \mathit{String})$.
\item[Conversation:] an instance of the Conversation Management ADT
  (\hyperref[mConv]{M2}).
\item[Status:] one of $\{\mathit{queued}$, $\mathit{running}$, $\mathit{completed}$, $\mathit{failed}$, $\mathit{cancelled}$, $\mathit{timedOut}\}$.
\item[Policy:] a tuple $(\mathit{table}: \mathit{String} \to \mathit{ModelName}$, $\mathit{threshold}: \mathbb{R})$, the task-to-model table and the complexity
  threshold.
\item[Admission:] one of $\{\mathit{prefill}$, $\mathit{queued}$, $\mathit{retry}\}$, the
  decision for a fallback prefill.
\item[Outcome:] one of $\{\mathit{reused}$, $\mathit{fellBack}$, $\mathit{failed}\}$, the
  result of one request for the fallback window.
\item[Entry:] a run-log record, a tuple $(\mathit{id}: \mathit{RequestID}$, $\mathit{source}: \mathit{ModelName}$, $\mathit{target}: \mathit{ModelName}$, $\mathit{status}: \mathit{Status}$, $\mathit{reused}: \mathbb{B}$, $\mathit{timings}:
  \text{Seq of } \mathbb{R}$, $\mathit{translator}: \mathit{String}$, $\mathit{checkpoint}: \mathit{String}$, $\mathit{translated}: \mathbb{B})$.
\end{description}

\subsubsection*{Evaluation and fitting}

\begin{description}
\item[Score:] a tuple $(\mathit{equivalent}: \mathbb{B}$, $\mathit{correct}: \mathbb{B}
  \cup \{\bot\})$; \emph{correct} is $\bot$ when the turn has no reference answer.
\item[WorkloadTurn:] a tuple $(\mathit{message}: \mathit{String}$, $\mathit{gold}:
  \mathit{String} \cup \{\bot\})$, one user message and its reference answer, or
  $\bot$ when it has none.
\item[Workload:] a tuple $(\mathit{name}: \mathit{String}$, $\mathit{conversations}:
  \text{Seq of (Seq of WorkloadTurn)}$, $\mathit{settings}: \mathit{Settings})$.
\item[Report:] Seq of $(\mathit{native}: \mathit{Entry}$, $\mathit{reused}:
  \mathit{Entry}$, $\mathit{score}: \mathit{Score})$, one row per request of a
  workload.
\item[Approval:] one of $\{\mathit{approved}$, $\mathit{notApproved}\}$.
\item[SetupRecord:] a tuple $(\mathit{pair}: \mathit{ModelName} \times \mathit{ModelName}$, $\mathit{revision}: \mathit{Revision}$, $\mathit{version}: \mathit{String})$.
\item[Evidence:] a tuple $(\mathit{heldOut}: \mathbb{R}$, $\mathit{split}: \mathit{String}$, $\mathit{outputPassed}: \mathbb{B}$, $\mathit{privacyPassed}: \mathbb{B}$, $\mathit{failurePassed}: \mathbb{B})$, the tests recorded for a setup.
\item[Corpus:] a set of paired conversations.
\item[Split:] a tuple $(\mathit{fitting}$, $\mathit{selection}$, $\mathit{test})$ of
  disjoint sets of conversation groups.
\item[Method:] the name of a fitting procedure, a String.
\end{description}

The access programs in each module are given by a table with columns \emph{In},
\emph{Out} and \emph{Exceptions}. Each input is written as \emph{name}: Type. An
entry of ``--'' means none. A program with several outputs returns a tuple,
written $T_1 \times T_2$, and $\bot$ means no value. A library keeps no state
between calls, but it may read or write environment variables. In a template module, the access programs other than
the constructor (\emph{new}) act on the instance they are called on. In an
abstract object, every access program except \emph{init} raises
\emph{NotInitialized} when it is called before \emph{init}.

\newpage
\section{Hardware-Hiding Module}

\subsection{Accelerator Interface (M1)}\label{mAccel}

\textbf{Module:} Accelerator Interface. \textbf{Type:} Library.

\textbf{Uses:} OS.

\subsubsection*{Syntax}

\begin{tabularx}{\textwidth}{l X l X}
\toprule
\textbf{Name} & \textbf{In} & \textbf{Out} & \textbf{Exceptions}\\
\midrule
allocate & \emph{device}: DeviceName, \emph{dtype}: DType & DeviceHandle & UnsupportedDevice, UnsupportedDType\\
\bottomrule
\end{tabularx}

\subsubsection*{State Variables}

None.

\subsubsection*{Environment Variables}

\begin{description}
\item[devices:] $\mathcal{P}$ DeviceName, the devices present on the machine.
\end{description}

\newpage
\section{Behaviour-Hiding Modules}

\subsection{Conversation Management (M2)}\label{mConv}

\textbf{Template Module:} Conversation Management. \textbf{Type:} Abstract Data Type. \textbf{Exported type:} Conversation.

\textbf{Uses:} \hyperref[mCtrl]{Request Control (M3)}.

\subsubsection*{Syntax}

\begin{tabularx}{\textwidth}{l X l X}
\toprule
\textbf{Name} & \textbf{In} & \textbf{Out} & \textbf{Exceptions}\\
\midrule
new Conversation & \emph{model}: ModelName, \emph{domain}: ReuseDomain $\cup$ \{$\bot$\} & Conversation & --\\
send & \emph{message}: String, \emph{maxNewTokens}: $\mathbb{N}$ & Answer & InvalidMessage, RequestFailed\\
\bottomrule
\end{tabularx}

\emph{send} both extends the history and returns the answer, because a turn is
one exchange.

\subsubsection*{State Variables}

\begin{description}
\item[history:] Seq of Turn, the user and assistant messages in order.
\item[cache:] Cache $\cup$ \{$\bot$\}, the cache after the last successful answer; $\bot$ before the first answer.
\item[domain:] ReuseDomain $\cup$ \{$\bot$\}, the reuse domain given by the host at creation.
\item[current:] ModelName, the model that produced \emph{cache}.
\item[text:] String, the conversation text that \emph{cache} represents.
\item[boundary:] $\mathbb{N}$, the character offset in \emph{text} where \emph{cache} ends.
\end{description}

\subsubsection*{Environment Variables}

None.

\newpage
\subsection{Request Control (M3)}\label{mCtrl}

\textbf{Module:} Request Control. \textbf{Type:} Library.

\textbf{Uses:} \hyperref[mSelect]{Model Selection (M5)}, \hyperref[mReuse]{Reuse Policy (M6)}, \hyperref[mHandoff]{Cache Handoff (M7)}, \hyperref[mCompat]{Cache Compatibility (M8)}, \hyperref[mModelRun]{Model Runtime (M12)}, \hyperref[mLife]{Request Lifecycle (M9)}, \hyperref[mLog]{Run Logging (M10)}.

\subsubsection*{Syntax}

\begin{tabularx}{\textwidth}{l X l X}
\toprule
\textbf{Name} & \textbf{In} & \textbf{Out} & \textbf{Exceptions}\\
\midrule
process & \emph{req}: Request, \emph{prior}: Cache $\cup$ \{$\bot$\}, \emph{boundary}: $\mathbb{N}$ & Answer $\times$ Cache & InvalidRequest(ErrorCode), RequestFailed, Overloaded, ContextTooLong\\
\bottomrule
\end{tabularx}

\subsubsection*{State Variables}

None.

\subsubsection*{Environment Variables}

None.

\newpage
\subsection{Comparison Benchmark (M4)}\label{mBench}

\textbf{Module:} Comparison Benchmark. \textbf{Type:} Library.

\textbf{Uses:} \hyperref[mCtrl]{Request Control (M3)}, \hyperref[mQuality]{Quality Measurement (M11)}, \hyperref[mLog]{Run Logging (M10)}.

\subsubsection*{Syntax}

\begin{tabularx}{\textwidth}{l X l X}
\toprule
\textbf{Name} & \textbf{In} & \textbf{Out} & \textbf{Exceptions}\\
\midrule
runWorkload & \emph{workload}: Workload & Report & MismatchedWorkload\\
\bottomrule
\end{tabularx}

\subsubsection*{State Variables}

None.

\subsubsection*{Environment Variables}

None.

\newpage
\subsection{Model Selection (M5)}\label{mSelect}

\textbf{Module:} Model Selection. \textbf{Type:} Abstract Object.

\textbf{Uses:} None.

\subsubsection*{Syntax}

\begin{tabularx}{\textwidth}{l X l X}
\toprule
\textbf{Name} & \textbf{In} & \textbf{Out} & \textbf{Exceptions}\\
\midrule
init & \emph{policy}: Policy, \emph{device}: DeviceHandle & -- & AlreadyInitialized, InvalidPolicy, ClassifierUnavailable\\
select & \emph{prompt}: String & ModelName & NotInitialized, EmptyPrompt, ClassificationFailed\\
\bottomrule
\end{tabularx}

\subsubsection*{State Variables}

\begin{description}
\item[policy:] Policy, the task-to-model table and the complexity threshold.
\item[classifier:] ClassifierModel, the loaded classifier.
\item[initialized:] $\mathbb{B}$, false until \emph{init} is called.
\end{description}

\subsubsection*{Environment Variables}

\begin{description}
\item[weights:] File, the classifier checkpoint on disk.
\end{description}

\newpage
\subsection{Reuse Policy (M6)}\label{mReuse}

\textbf{Module:} Reuse Policy. \textbf{Type:} Abstract Object.

\textbf{Uses:} \hyperref[mMap]{Cache Mapper (M14)}, \hyperref[mLimits]{Limits (M16)}.

\subsubsection*{Syntax}

\begin{tabularx}{\textwidth}{l X l X}
\toprule
\textbf{Name} & \textbf{In} & \textbf{Out} & \textbf{Exceptions}\\
\midrule
init & -- & -- & AlreadyInitialized\\
qualify & \emph{setup}: SetupRecord, \emph{evidence}: Evidence & -- & NotInitialized, IncompleteEvidence\\
allowed & \emph{target}: ModelName, \emph{req}: Request & $\mathbb{B}$ & NotInitialized\\
\bottomrule
\end{tabularx}

\subsubsection*{State Variables}

\begin{description}
\item[approved:] $\mathcal{P}$ SetupRecord, the setups approved for reuse.
\item[initialized:] $\mathbb{B}$, false until \emph{init} is called.
\end{description}

\subsubsection*{Environment Variables}

None.

\newpage
\subsection{Cache Handoff (M7)}\label{mHandoff}

\textbf{Module:} Cache Handoff. \textbf{Type:} Library.

\textbf{Uses:} \hyperref[mMap]{Cache Mapper (M14)}, \hyperref[mCompat]{Cache Compatibility (M8)}, \hyperref[mLife]{Request Lifecycle (M9)}, \hyperref[mModelRun]{Model Runtime (M12)}.

\subsubsection*{Syntax}

\begin{tabularx}{\textwidth}{l X l X}
\toprule
\textbf{Name} & \textbf{In} & \textbf{Out} & \textbf{Exceptions}\\
\midrule
handoff & \emph{cache}: Cache, \emph{target}: ModelName, \emph{conversation}: String, \emph{boundary}: $\mathbb{N}$, \emph{maxNewTokens}: $\mathbb{N}$ & Answer $\times$ Cache & UnsupportedDirection, HandoffFailed\\
\bottomrule
\end{tabularx}

\subsubsection*{State Variables}

None.

\subsubsection*{Environment Variables}

None.

\newpage
\subsection{Cache Compatibility (M8)}\label{mCompat}

\textbf{Module:} Cache Compatibility. \textbf{Type:} Library.

\textbf{Uses:} \hyperref[mCache]{Cache Store (M13)}, \hyperref[mModelRun]{Model Runtime (M12)}, \hyperref[mLimits]{Limits (M16)}.

\subsubsection*{Syntax}

\begin{tabularx}{\textwidth}{l X l X}
\toprule
\textbf{Name} & \textbf{In} & \textbf{Out} & \textbf{Exceptions}\\
\midrule
seal & \emph{cache}: Cache, \emph{checkpoint}: CheckpointRef, \emph{targetRevision}: Revision & CacheRecord & --\\
check & \emph{cache}: Cache, \emph{record}: CacheRecord $\cup$ \{$\bot$\}, \emph{target}: ModelName, \emph{conversation}: String, \emph{boundary}: $\mathbb{N}$ & Verdict & --\\
canary & \emph{predicted}: Cache & $\mathbb{B}$ & CanaryUnavailable\\
\bottomrule
\end{tabularx}

\subsubsection*{State Variables}

None.

\subsubsection*{Environment Variables}

None.

\newpage
\subsection{Request Lifecycle (M9)}\label{mLife}

\textbf{Module:} Request Lifecycle. \textbf{Type:} Abstract Object.

\textbf{Uses:} \hyperref[mCache]{Cache Store (M13)}, \hyperref[mModelRun]{Model Runtime (M12)}, \hyperref[mLimits]{Limits (M16)}.

\subsubsection*{Syntax}

\begin{tabularx}{\textwidth}{l X l X}
\toprule
\textbf{Name} & \textbf{In} & \textbf{Out} & \textbf{Exceptions}\\
\midrule
init & -- & -- & AlreadyInitialized\\
start & \emph{id}: RequestID & -- & NotInitialized, DuplicateRequest\\
finish & \emph{id}: RequestID, \emph{status}: Status & -- & NotInitialized, UnknownRequest\\
cancel & \emph{id}: RequestID & -- & NotInitialized, UnknownRequest\\
status & \emph{id}: RequestID & Status & NotInitialized, UnknownRequest\\
admit & \emph{id}: RequestID & Admission & NotInitialized, UnknownRequest\\
recordOutcome & \emph{id}: RequestID, \emph{outcome}: Outcome & -- & NotInitialized, UnknownRequest\\
release & \emph{cache}: Cache & -- & NotInitialized\\
holdPredicted & \emph{cache}: Cache & -- & NotInitialized\\
markUsed & \emph{cache}: Cache & -- & NotInitialized\\
expireUnused & -- & -- & NotInitialized\\
reuseEnabled & -- & $\mathbb{B}$ & NotInitialized\\
\bottomrule
\end{tabularx}

\emph{admit} both decides and takes the queue slot, so that two requests cannot
see the same free slot.

\subsubsection*{State Variables}

\begin{description}
\item[states:] RequestID $\to$ Status, the current state of each request.
\item[window:] Seq of Outcome, the outcomes of the last \texttt{FALLBACK\_WINDOW} requests.
\item[queued:] $\mathbb{N}$, the number of fallback prefills waiting, at most \texttt{QUEUE\_LIMIT}.
\item[enabled:] $\mathbb{B}$, whether cross-model reuse is on; false once the fallback rate exceeds \texttt{FALLBACK\_RATE\_LIMIT}.
\item[unused:] $\mathcal{P}$(Cache $\times$ $\mathbb{R}$), the predicted caches not yet used, each with the time it was made. Their total size is at most \texttt{PREDICTED\_CACHE\_LIMIT}, and each expires after \texttt{PREDICTED\_CACHE\_TTL} (SR-14).
\item[initialized:] $\mathbb{B}$, false until \emph{init} is called.
\end{description}

\subsubsection*{Environment Variables}

\begin{description}
\item[clock:] $\mathbb{R}$, the time used for timeouts and expiry.
\end{description}

\newpage
\subsection{Run Logging (M10)}\label{mLog}

\textbf{Module:} Run Logging. \textbf{Type:} Library.

\textbf{Uses:} None.

\subsubsection*{Syntax}

\begin{tabularx}{\textwidth}{l X l X}
\toprule
\textbf{Name} & \textbf{In} & \textbf{Out} & \textbf{Exceptions}\\
\midrule
record & \emph{entry}: Entry & -- & LogWriteFailed\\
\bottomrule
\end{tabularx}

\subsubsection*{State Variables}

None.

\subsubsection*{Environment Variables}

\begin{description}
\item[logFile:] File, the destination for records.
\end{description}

\newpage
\subsection{Quality Measurement (M11)}\label{mQuality}

\textbf{Module:} Quality Measurement. \textbf{Type:} Library.

\textbf{Uses:} None.

\subsubsection*{Syntax}

\begin{tabularx}{\textwidth}{l X l X}
\toprule
\textbf{Name} & \textbf{In} & \textbf{Out} & \textbf{Exceptions}\\
\midrule
score & \emph{answer}: Answer, \emph{native}: Answer, \emph{gold}: String $\cup$ \{$\bot$\} & Score & EmptyReference\\
taskScore & \emph{answer}: Answer, \emph{gold}: String & $\mathbb{R}$ & EmptyReference\\
\bottomrule
\end{tabularx}

\subsubsection*{State Variables}

None.

\subsubsection*{Environment Variables}

None.

\newpage
\section{Software Decision Modules}

\subsection{Model Runtime (M12)}\label{mModelRun}

\textbf{Module:} Model Runtime. \textbf{Type:} Abstract Object.

\textbf{Uses:} \hyperref[mAccel]{Accelerator Interface (M1)}.

\subsubsection*{Syntax}

\begin{tabularx}{\textwidth}{l X l X}
\toprule
\textbf{Name} & \textbf{In} & \textbf{Out} & \textbf{Exceptions}\\
\midrule
init & -- & -- & AlreadyInitialized\\
load & \emph{model}: ModelName, \emph{revision}: Revision, \emph{device}: DeviceHandle & -- & NotInitialized, LoadFailed\\
tokenize & \emph{model}: ModelName, \emph{text}: String & Seq of Token & NotInitialized, ModelNotLoaded\\
revision & \emph{model}: ModelName & Revision & NotInitialized, ModelNotLoaded\\
generate & \emph{model}: ModelName, \emph{tokens}: Seq of Token, \emph{prior}: Cache $\cup$ \{$\bot$\}, \emph{maxNewTokens}: $\mathbb{N}$ & Answer $\times$ Cache & NotInitialized, ModelNotLoaded, ContextTooLong\\
\bottomrule
\end{tabularx}

\subsubsection*{State Variables}

\begin{description}
\item[loaded:] ModelName $\to$ Model, the loaded checkpoints.
\item[initialized:] $\mathbb{B}$, false until \emph{init} is called.
\end{description}

\subsubsection*{Environment Variables}

\begin{description}
\item[weights:] File, the checkpoint files on disk.
\end{description}

\newpage
\subsection{Cache Store (M13)}\label{mCache}

\textbf{Template Module:} Cache Store. \textbf{Type:} Abstract Data Type. \textbf{Exported type:} Cache.

\textbf{Uses:} None.

\subsubsection*{Syntax}

\begin{tabularx}{\textwidth}{l X l X}
\toprule
\textbf{Name} & \textbf{In} & \textbf{Out} & \textbf{Exceptions}\\
\midrule
new Cache & \emph{model}: ModelName, \emph{tokens}: Seq of Token, \emph{data}: KVData, \emph{complete}: $\mathbb{B}$ & Cache & --\\
length & -- & $\mathbb{N}$ & --\\
provenance & -- & Provenance & --\\
kv & -- & KVData & --\\
copy & -- & Cache & --\\
\bottomrule
\end{tabularx}

\subsubsection*{State Variables}

\begin{description}
\item[data:] KVData, the keys and values of the cache.
\item[model:] ModelName, the model that produced the cache.
\item[tokens:] Seq of Token, the exact token IDs represented.
\item[complete:] $\mathbb{B}$, true when the last generated token ended the turn.
\end{description}

\subsubsection*{Environment Variables}

None.

\newpage
\subsection{Cache Mapper (M14)}\label{mMap}

\textbf{Module:} Cache Mapper. \textbf{Type:} Abstract Object.

\textbf{Uses:} \hyperref[mCache]{Cache Store (M13)}, \hyperref[mModelRun]{Model Runtime (M12)}.

\subsubsection*{Syntax}

\begin{tabularx}{\textwidth}{l X l X}
\toprule
\textbf{Name} & \textbf{In} & \textbf{Out} & \textbf{Exceptions}\\
\midrule
init & -- & -- & AlreadyInitialized\\
register & \emph{source}: ModelName, \emph{target}: ModelName, \emph{checkpoint}: CheckpointRef & -- & NotInitialized, DuplicateDirection, InvalidCheckpoint\\
lookup & \emph{source}: ModelName, \emph{target}: ModelName & CheckpointRef & NotInitialized, UnsupportedDirection, UnavailableCheckpoint\\
translate & \emph{cache}: Cache, \emph{target}: ModelName, \emph{conversation}: String, \emph{boundary}: $\mathbb{N}$ & Cache $\times$ Seq of Token & NotInitialized, UnsupportedDirection, InvalidCheckpoint\\
\bottomrule
\end{tabularx}

\subsubsection*{State Variables}

\begin{description}
\item[entries:] (ModelName $\times$ ModelName) $\to$ CheckpointRef, the registered checkpoint for each direction.
\item[maps:] (ModelName $\times$ ModelName) $\to$ Map, the loaded map for each registered direction.
\item[initialized:] $\mathbb{B}$, false until \emph{init} is called.
\end{description}

\subsubsection*{Environment Variables}

\begin{description}
\item[checkpoints:] File, the translator checkpoint files on disk.
\end{description}

\newpage
\subsection{Predictor Fitting (M15)}\label{mFit}

\textbf{Module:} Predictor Fitting. \textbf{Type:} Library. \textbf{Runs:} offline only.

\textbf{Uses:} \hyperref[mModelRun]{Model Runtime (M12)}, \hyperref[mCache]{Cache Store (M13)}.

\subsubsection*{Syntax}

\begin{tabularx}{\textwidth}{l X l X}
\toprule
\textbf{Name} & \textbf{In} & \textbf{Out} & \textbf{Exceptions}\\
\midrule
split & \emph{corpus}: Corpus, \emph{selectionGroups}: $\mathbb{N}$, \emph{testGroups}: $\mathbb{N}$ & Split & InsufficientData\\
fit & \emph{split}: Split, \emph{method}: Method & CheckpointRef & OverlappingSplit, FitFailed\\
\bottomrule
\end{tabularx}

\subsubsection*{State Variables}

None.

\subsubsection*{Environment Variables}

\begin{description}
\item[checkpoints:] File, the fitted predictor checkpoints that \emph{fit} writes.
\end{description}

\newpage
\subsection{Limits (M16)}\label{mLimits}

\textbf{Module:} Limits. \textbf{Type:} Library.

\textbf{Uses:} None.

\subsubsection*{Exported Constants}

The values come from the
VnV Plan \citep[Appendix, Symbolic Parameters]{TandemVnV}
and the SRS; those marked provisional are confirmed on development data (SRS OI3).

\begin{center}
\renewcommand{\arraystretch}{1.2}
\begin{tabular}{l l l}
\toprule
\textbf{Constant} & \textbf{Value} & \textbf{Source}\\
\midrule
\texttt{SCORE\_ATOL} & 0.001 & SRS\\
\texttt{SCORE\_RTOL} & 0.001 & SRS\\
\texttt{GREEDY\_TOKENS} & 32 & SRS\\
\texttt{REGRESSION\_LIMIT} & 0.05 & SRS NFR-07, NFR-08\\
\texttt{QUALITY\_LOSS\_LIMIT} & 0.05 (answer F1, absolute) & VnV; HA SR-09 (provisional)\\
\texttt{CANARY\_TOKENS} & 16 & VnV; HA SR-10 (provisional)\\
\texttt{CANARY\_LIMIT} & 0.5 (relative K/V error) & VnV; HA SR-10 (provisional)\\
\texttt{FALLBACK\_WINDOW} & 100 requests & VnV; HA SR-13 (provisional)\\
\texttt{FALLBACK\_RATE\_LIMIT} & 0.5 & VnV; HA SR-13 (provisional)\\
\texttt{QUEUE\_LIMIT} & 32 requests & VnV; HA SR-13 (provisional)\\
\texttt{PREDICTED\_CACHE\_LIMIT} & 2 GiB & VnV; HA SR-14 (provisional)\\
\texttt{PREDICTED\_CACHE\_TTL} & 60 s & VnV; HA SR-14 (provisional)\\
\bottomrule
\end{tabular}
\end{center}

\subsubsection*{Syntax}

None. Limits exports constants only.

\subsubsection*{State Variables}

None.

\subsubsection*{Environment Variables}

None.

\newpage
\section*{Appendix --- Generative AI Use Disclosure}

\subsection*{AI Tools Used}

\begin{itemize}
\item Claude Code, a command-line tool, using the model Claude Opus 5.5
  (model ID \texttt{claude-opus-5-5}).
\end{itemize}
No other AI tools were used in drafting this document.

\subsection*{How and Where Used}

\begin{itemize}
\item \textbf{Reviewing the document.} The draft was checked against the course
  MIS checklist, the template and the Rev $-1$ rubric. The tool listed types that
  were used but not defined, abstract objects without initialization, arguments
  whose meaning was unclear, exceptions that could not occur, and uses that
  disagreed with the Module Guide or the code.
\item \textbf{Modifying text.} The tool removed sections that Rev $-1$ does not
  require, removed the eleven modules merged in the Module Guide (27 to 16
  modules), and rewrote the derived types and the access-program tables of all
  16 modules.
\end{itemize}

\subsection*{Verification of AI-Generated Content}

\begin{itemize}
\item Every change was proposed as a plan and approved by a team member before
  it was made.
\item Argument names and uses were checked against the source code in
  \texttt{src/}, and module names, types and uses against the Module Guide.
\item The document was rebuilt with no LaTeX errors or undefined references.
\end{itemize}
The team reviewed the final text and is responsible for its accuracy.

\bibliographystyle{plainnat}
\bibliography{../../../refs/References}

\end{document}
